\chapter{Căn Bệnh Con Nhà Người Ta}
\markboth{Căn Bệnh Con Nhà Người Ta}{Căn Bệnh Con Nhà Người Ta}

\section{Một Câu So Sánh, Mười Vết Xước}
\begin{truyen}{Một Câu So Sánh, Mười Vết Xước}{One Comparison, Ten Scratches}
\chuhoa{H}{ọc sinh:} ``Ba mẹ em cứ lôi con nhà người ta ra. Nghe riết em không biết mình đang sống đời mình hay đời người ta nữa.''
\textit{(Student: ``My parents keep bringing up other people's children. After hearing it so often, I no longer know whether I am living my own life or theirs.'')}

\teacherpair{So sánh là cách dạy rất rẻ nhưng hậu quả rất đắt. Nó có thể đẩy một đứa trẻ chạy nhanh hơn trong chốc lát, nhưng cũng rất dễ làm nó xấu hổ về chính mình, ganh ghét người khác, và sống trong cảm giác luôn thiếu.}{``Comparison is a cheap teaching method with expensive consequences. It may push a child to run faster for a while, but it also easily makes that child ashamed of themselves, resentful toward others, and permanently inadequate.''}
\end{truyen}

\section{Con Nhà Người Ta Cũng Có Một Cơn Khổ Riêng}
\begin{truyen}{Con Nhà Người Ta Cũng Có Một Cơn Khổ Riêng}{Other People's Children Suffer Too}
\chuhoa{H}{ọc sinh:} ``Nhưng đúng là có mấy bạn giỏi thật mà thầy.''
\textit{(Student: ``But some of those other children are genuinely excellent, teacher.'')}

\teacherpair{Đúng. Nhưng người lớn hay chỉ nhìn phần sáng của con nhà người ta và quên rằng nhà nào cũng có góc tối. Có bạn học giỏi nhưng áp lực, có bạn ngoan nhưng cô độc, có bạn thành tích đẹp mà sống trong nỗi sợ làm bố mẹ thất vọng. Lấy một lát cắt đẹp của người khác để đánh vào cả con người của con mình là rất tàn nhẫn.}{``True. But adults usually look only at the bright side of other people's children and forget that every home has shadowed corners. Some students do well but live under pressure, some behave well but feel lonely, and some have beautiful results while living in fear of disappointing their parents. Using one polished slice of another child to attack your own child is cruel.''}
\end{truyen}

\section{Muốn Con Tốt Hơn, Đừng Xóa Giá Trị Của Nó}
\begin{truyen}{Muốn Con Tốt Hơn, Đừng Xóa Giá Trị Của Nó}{If You Want a Child to Improve, Do Not Erase Their Worth}
\chuhoa{H}{ọc sinh:} ``Bị so riết em không muốn cố nữa.''
\textit{(Student: ``After being compared for so long, I do not even want to try anymore.'')}

\teacherpair{Đó là phản ứng rất thường. Một người bị phủ nhận liên tục sẽ hoặc gồng lên, hoặc gục xuống. Muốn con tốt hơn thì phải chỉ ra cái cần sửa mà vẫn giữ được cảm giác rằng nó có giá trị, có thể lớn lên, và không phải bản lỗi của một mẫu người khác.}{``That is a common reaction. A person who is constantly invalidated either hardens or collapses. If you want a child to improve, you must point out what needs changing while preserving their sense of worth, their ability to grow, and the truth that they are not a faulty copy of someone else.''}
\end{truyen}

\begin{baihoc}
So sánh có thể tạo chuyển động, nhưng rất ít khi tạo trưởng thành lành mạnh.
\textit{(Comparison may create movement, but it rarely creates healthy growth.)}
\end{baihoc}

\begin{ghinhoanh}
\textbf{Short English to remember:}

Comparison can motivate.

It can also wound identity.
\end{ghinhoanh}

\begin{tuhocnhanh}
1. Câu so sánh nào từng làm em đau lâu nhất? \textit{(Which comparison hurt you for the longest time?)}

2. Em có đang vô thức dùng người khác để chà lên giá trị của chính mình không? \textit{(Are you unconsciously using others to grind down your own worth?)}
\end{tuhocnhanh}

\ngancach
